% 2007 年湖南卷文科第 7 题：河流水位的频率分布直方图。
% 来源 PDF：普通高考/2007/2007湖南文.pdf，SHA256 8e889ae321ffa864b24c505b40e7e96956b301487f5a106264b7fcd9f943e068
% 原页：第 1 页（A3 三栏）第 2 栏上部，印刷页 381，图在题干下方。
% 重绘依据：横轴水位 30–51 米、组距 1 米；纵轴为频率/组距，标出 0.5%、1%、2% 三条虚线。
%   逐柱高度按原图柱顶量取（以 1% 线为标尺），再按「各组频率之和为 100%」归一：
%   30–31 为 2.6%、31–32 为 3.1%、32–33 为 3.4%，其后渐升，39–40 达峰 8.0%，
%   43–47 有起伏；尾三组按原图标注取 48–49 为 2.0%、49–50 为 1.0%、50–51 为 0.5%。
%   原图只在两端标出 30–33 与 48–51，此处照录。
\begin{tikzpicture}[scale=0.60, semithick, every node/.style={inner sep=1.5pt}]
  % 坐标轴
  \draw[->] (29.7,0) -- (51.9,0) node[anchor=west] {水位（米）};
  \draw[->] (30,0) -- (30,7.6);
  \node[align=center, anchor=south] at (30,7.65) {频率\\[-1pt]组距};
  % 断裂记号
  \draw (29.85,0) -- (29.95,0.12) -- (30.05,-0.02);
  % 三条水平虚线
  \draw[dashed] (30,0.5) -- (51.4,0.5);
  \node[anchor=east] at (29.92,0.5) {$0.5\%$};
  \draw[dashed] (30,1) -- (51.4,1);
  \node[anchor=east] at (29.92,1) {$1\%$};
  \draw[dashed] (30,2) -- (51.4,2);
  \node[anchor=east] at (29.92,2) {$2\%$};
  % 各柱
  \foreach \x/\h in {
    30/2.6, 31/3.1, 32/3.4, 33/3.9, 34/4.5, 35/5.1, 36/5.8,
    37/6.2, 38/6.6, 39/8.0, 40/7.5, 41/6.9, 42/6.3, 43/5.2,
    44/5.8, 45/5.4, 46/5.8, 47/4.8, 48/2.0, 49/1.0, 50/0.5} {
    \draw (\x,\h) rectangle (\x+1,0);
  }
  % 横轴刻度标签（照原图只标两端）
  \foreach \x in {30,31,32,33,48,49,50,51} {
    \node[anchor=north] at (\x,-0.04) {\x};
  }
\end{tikzpicture}
